% from ega0-23-fr.tex
\oldpage[0\textsubscript{IV-1}]{218}

le nilradical de $B$, le noyau de $v$ est $K'\otimes_A\mathfrak N$, donc il est
contenu dans le nilradical de $K'\otimes_AB$~; on en conclut aussitôt que si
l'image par $h=w\circ v\circ u$ d'un élément $x\in A'\otimes_AB$ est
nilpotent, $x$ est nilpotent, ce qui prouve le lemme.
\end{proof}

\begin{proposition}[23.2.4]
\label{0.23.2.4-fr}
Soient $A$ un anneau intègre, $K$ son corps des fractions, $B$ un anneau
noethérien, $\mathfrak q_i$ ($1\leqslant i\leqslant n$) ses idéaux premiers
minimaux, $\varphi:A\to B$ un homomorphisme d'anneaux. On suppose que les
$B/\mathfrak q_i$ sont des anneaux japonais, et que $B$ soit un $A$-module
plat. Soit $K'$ une extension finie de $K$, et soit $(C_\lambda)$ la famille
filtrante croissante des sous-anneaux de $K'$ qui sont des $A$-algèbres
finies et admettent $K'$ pour corps des fractions~; la réunion $A'$ des
$C_\lambda$ est donc la fermeture intégrale de $A$ dans $K'$. Alors~:
\begin{enumerate}
  \item[(i)] Il existe un indice $\alpha$ tel que pour $\lambda\geqslant
  \alpha$, l'homomorphisme canonique
  \begin{equation}
  \label{0.23.2.4.1-fr}
    (C_\alpha\otimes_AB)_{\mathrm{red}}
    \longrightarrow(C_\lambda\otimes_AB)_{\mathrm{red}}
    \tag{23.2.4.1}
  \end{equation}
  soit bijectif.
  \item[(ii)] Si en outre $B$ est un $A$-module fidèlement plat, le morphisme
  $\operatorname{Spec}(A')\to\operatorname{Spec}(C_\alpha)$ est radiciel.
\end{enumerate}
\end{proposition}

\begin{proof}
\emph{(i)} Soit $B_1=B_{\mathrm{red}}$, et soit $R$ l'anneau total des
fractions de $B_1$~; comme $A$ est intègre et $B$ un $A$-module plat, tout
élément $\neq0$ de $A$ est $B$-régulier (0\textsubscript{I}, 6.3.4), donc
l'homomorphisme composé $A\xrightarrow{\varphi}B\to B_1$ se prolonge en un
homomorphisme $K\to R$~; on peut alors écrire
$K'\otimes_AR=(K'\otimes_AK)\otimes_KR$, et comme $K'\otimes_AK$
s'identifie à $K'$, on a $K'\otimes_AR=K'\otimes_KR$. Comme $R$ est composé
direct des corps des fractions $L_i$ des $B/\mathfrak q_i$,
$K'\otimes_AR$ est composé direct des $K'\otimes_KL_i$, et par suite
$(K'\otimes_AR)_{\mathrm{red}}$ est composé direct d'un nombre fini
d'extensions finies des $L_i$. En vertu de l'hypothèse, la fermeture
intégrale de $B/\mathfrak q_i$ dans une extension finie de $L_i$ est un
$(B/\mathfrak q_i)$-module de type fini, donc un $B$-module de type fini~;
on en conclut que la fermeture intégrale $B'$ de $B$ dans
$(K'\otimes_AR)_{\mathrm{red}}$ est un $B$-module de type fini. En vertu de
(23.2.3), les $(C_\lambda\otimes_AB)_{\mathrm{red}}$ s'identifient
canoniquement à des sous-anneaux de $(K'\otimes_AR)_{\mathrm{red}}$ qui sont
des $B$-algèbres finies, donc contenues dans $B'$. Comme $B$ est noethérien
et $B'$ un $B$-module de type fini, la famille filtrante des
$(C_\lambda\otimes_AB)_{\mathrm{red}}$ admet un \emph{plus grand élément}
$(C_\alpha\otimes_AB)_{\mathrm{red}}$, ce qui prouve (i).

\emph{(ii)} Comme $B$ est un $A$-module fidèlement plat, il suffit
(IV, 2.6.1, (v)) de montrer que le morphisme
$\operatorname{Spec}(A'\otimes_AB)\to
\operatorname{Spec}(C_\alpha\otimes_AB)$ est radiciel, ou, ce qui revient au
même (I, 5.1.6), que le morphisme
$\operatorname{Spec}(A'\otimes_AB)_{\mathrm{red}}\to
\operatorname{Spec}(C_\alpha\otimes_AB)_{\mathrm{red}}$ est radiciel~; mais
on a $A'\otimes_AB=\varinjlim_\lambda(C_\lambda\otimes_AB)$, donc
(IV, 5.13.2)
$(A'\otimes_AB)_{\mathrm{red}}=
\varinjlim_\lambda(C_\lambda\otimes_AB)_{\mathrm{red}}$, et la conclusion
résulte de (i).
\end{proof}

\begin{corollary}[23.2.5]
\label{0.23.2.5-fr}
Soient $A$ un anneau local noethérien intègre, $K$ son corps des fractions,
$K'$ une extension finie de $K$. Soit $(C_\lambda)$ la famille filtrante
croissante des sous-anneaux de $K'$ qui sont des $A$-algèbres finies (donc
des anneaux noethériens semi-locaux (Bourbaki, \emph{Alg. comm.}, chap.~IV,
§~2, n\textsuperscript{o}~5, cor.~3 de la prop.~9)) et admettent $K'$ pour
corps des fractions. Alors il existe $\alpha$ tel que l'homomorphisme
$(\widehat C_\alpha)_{\mathrm{red}}\to
(\widehat C_\lambda)_{\mathrm{red}}$ soit un isomorphisme pour
$\lambda\geqslant\alpha$, et si $A'$ est la fermeture intégrale de $A$ dans
$K'$, le morphisme
$\operatorname{Spec}(A')\to\operatorname{Spec}(C_\alpha)$ est radiciel
(cf. (23.2.2)).
\end{corollary}

